%% toc: P15 159 159.3,159.4,159.5,159.6
\section{Restricted integrals}
\label{sec:annex-p15-w02c-p15-007-restricted-integrals}

Throughout, \((\Omega,\mathcal F_0,P)\) is a probability space,
\(\mathcal G\subseteq\mathcal F_0\) is a sigma-algebra and
\(P|_{\mathcal G}\) is the restriction of
\cref{def:p15-w02c-p15-007-restricted-probability}. Integrals with respect
to \(P\) are taken on \((\Omega,\mathcal F_0,P)\), and integrals with
respect to \(P|_{\mathcal G}\) on \((\Omega,\mathcal G,P|_{\mathcal G})\),
each under \cref{def:p15-simple-integral},
\cref{def:p15-nonnegative-integral} and
\cref{def:p15-integrable-functions}.

\begin{annexitem}{annex:p15-w02c-p15-007-restricted}{thm:p15-w02c-p15-007-restricted-integral}{Approximation in \(\mathcal G\)}
\begin{proof}
\paragraph{Simple functions.}
Let \(s=\sum_{j=1}^m a_j\mathbf1_{A_j}\), where the \(A_j\in\mathcal G\)
are pairwise disjoint with union \(\Omega\) and the \(a_j\) are finite and
nonnegative. This is a representation of \(s\) as a simple function on
\((\Omega,\mathcal G)\). Since each \(A_j\) also lies in \(\mathcal F_0\),
the same sum is a representation of \(s\) as a simple function on
\((\Omega,\mathcal F_0)\). By \cref{def:p15-simple-integral}, and the
independence of the representation in
\cref{thm:p15-simple-integral-laws},
\[
 \int s\,d(P|_{\mathcal G})
   =\sum_{j=1}^m a_jP|_{\mathcal G}(A_j)
   =\sum_{j=1}^m a_jP(A_j)
   =\int s\,dP .
\]

\paragraph{Item 1.}
Let \(W\geq0\) be \(\mathcal G\)-measurable. On the probability space
\((\Omega,\mathcal G,P|_{\mathcal G})\), \(W\) is a nonnegative extended
real random variable. \Cref{thm:p15-measurable-limits}, applied on that
space, gives nonnegative simple functions \(s_n\) on
\((\Omega,\mathcal G)\) with \(s_n\leq s_{n+1}\) and \(s_n\to W\) at every
outcome; explicitly, the level sets
\(\{k2^{-n}\leq W<(k+1)2^{-n}\}\) and \(\{W\geq n\}\) in its formula are
inverse images under \(W\) of Borel sets and so belong to \(\mathcal G\).
Each \(s_n\) is therefore a simple function on both measurable spaces, and
the preceding paragraph gives
\(\int s_n\,d(P|_{\mathcal G})=\int s_n\,dP\) for every \(n\).
\Cref{thm:p15-monotone-integral} holds on every measure space; apply it
once on \((\Omega,\mathcal G,P|_{\mathcal G})\) and once on
\((\Omega,\mathcal F_0,P)\) to the same increasing sequence:
\[
 \int W\,d(P|_{\mathcal G})
   =\lim_{n\to\infty}\int s_n\,d(P|_{\mathcal G})
   =\lim_{n\to\infty}\int s_n\,dP
   =\int W\,dP .
\]
The common value may be \(\infty\). For \(A\in\mathcal G\), the function
\(W\mathbf1_A\) is nonnegative and \(\mathcal G\)-measurable: for a Borel
set \(B\) of extended reals, \(\{W\mathbf1_A\in B\}\) equals
\(\{W\in B\}\cap A\) if \(0\notin B\) and
\((\{W\in B\}\cap A)\cup A^c\) if \(0\in B\), and both lie in
\(\mathcal G\). Applying the display to \(W\mathbf1_A\) gives
\(\int_AW\,d(P|_{\mathcal G})=\int_AW\,dP\), by the definition of an
integral over an event on each space.

\paragraph{Item 2.}
For \(\mathcal G\)-measurable \(W\), the functions
\(W^+=\max(W,0)\) and \(W^-=\max(-W,0)\) are \(\mathcal G\)-measurable,
because \(\{W^+>c\}\) is \(\{W>c\}\) for \(c\geq0\) and \(\Omega\) for
\(c<0\), while \(\{W^->c\}\) is \(\{W<-c\}\) for \(c\geq0\) and \(\Omega\)
for \(c<0\);
rays of this kind generate the Borel sets of the extended reals. The sum
\(|W|=W^++W^-\) is nonnegative, and
\(\{|W|>c\}=\{W>c\}\cup\{W<-c\}\) for \(c\geq0\), so \(|W|\) is
\(\mathcal G\)-measurable. Item~1 with \(A=\Omega\) gives
\(\int|W|\,d(P|_{\mathcal G})=\int|W|\,dP=E|W|\), so one integral is
finite exactly when the other is.

For \(A\in\mathcal G\), \((W\mathbf1_A)^\pm=W^\pm\mathbf1_A\). Item~1,
applied to \(W^+\) and to \(W^-\), gives
\[
 \int_AW^+\,d(P|_{\mathcal G})=\int_AW^+\,dP,\qquad
 \int_AW^-\,d(P|_{\mathcal G})=\int_AW^-\,dP .
\]
By \cref{def:p15-integrable-functions}, the signed integral over \(A\)
exists for a measure exactly when the two numbers in its row are not both
infinite, and it then equals their difference. The two rows consist of the
same pair of numbers. Hence existence holds under \(P|_{\mathcal G}\)
exactly when it holds under \(P\), and the two signed integrals are the same
difference. The integrable case is included.

\paragraph{Item 3.}
Let \(V,W\) be \(\mathcal G\)-measurable extended real functions. Between
two distinct extended reals lies a rational number, so
\[
 \{V>W\}=\bigcup_{q\in\mathbb Q}\bigl(\{V>q\}\cap\{W<q\}\bigr).
\]
Each set on the right is an inverse image of a ray under a
\(\mathcal G\)-measurable function, and the union is countable, so
\(\{V>W\}\in\mathcal G\).

Suppose \(V\leq W\) holds \(P|_{\mathcal G}\)-almost everywhere. Then
\(\{V>W\}\subseteq E\) for some \(E\in\mathcal G\) with
\(P|_{\mathcal G}(E)=0\). Since \(E\in\mathcal F_0\) and \(P(E)=0\),
\(V\leq W\) holds \(P\)-almost everywhere. Conversely, suppose
\(\{V>W\}\subseteq E'\) for some \(E'\in\mathcal F_0\) with \(P(E')=0\).
Monotonicity of \(P\) gives \(P(\{V>W\})=0\). Because
\(\{V>W\}\in\mathcal G\), it is itself a \(P|_{\mathcal G}\)-null member
of \(\mathcal G\), outside which \(V\leq W\). For equality apply this to
\(\{V\neq W\}=\{V>W\}\cup\{W>V\}\), a union of two members of
\(\mathcal G\).

\paragraph{Second derivation of item 1 and of the real integrable case.}
Let \(\iota:\Omega\to\Omega\) be the identity. For \(A\in\mathcal G\),
\(\iota^{-1}(A)=A\in\mathcal F_0\), so \(\iota\) is a measurable map from
\((\Omega,\mathcal F_0)\) to \((\Omega,\mathcal G)\), that is, a random
element with values in \((\Omega,\mathcal G)\). By
\cref{def:p15-pushforward-law}, its law is
\(P_\iota(A)=P(\iota^{-1}(A))=P(A)=P|_{\mathcal G}(A)\). For \(A\in\mathcal G\)
and \(h=W\mathbf1_A\), \(h\) is \(\mathcal G\)-measurable and
\(h(\iota)=h\). If \(W\geq0\), or if \(W\) is real valued with
\(E|W|<\infty\), \cref{thm:p15-expectation-law-integral} gives
\[
 \int_AW\,dP=E[h(\iota)]=\int_\Omega h\,dP_\iota
            =\int_AW\,d(P|_{\mathcal G}).
\]
This agrees with items 1 and 2. The direct argument above also covers
extended real signed \(W\), which that theorem does not state.
\end{proof}
\end{annexitem}

\begin{annexitem}{annex:p15-w02c-p15-007-density}{thm:p15-w02c-p15-007-density-version}{The construction against \(P\)}
\begin{proof}
\paragraph{Item 1.}
By \cref{thm:p15-nonnegative-integral-laws}, item~4, the set functions
\(A\mapsto\int_AX^\pm\,dP\) are measures on \(\mathcal F_0\). Their values
on \(\mathcal G\) are \(\nu_\pm\). They vanish on \(\emptyset\), and a
countable disjoint family of members of \(\mathcal G\) is a countable
disjoint family of members of \(\mathcal F_0\); so \(\nu_\pm\) are
measures on \((\Omega,\mathcal G)\). Since \(0\leq X^\pm\leq|X|\),
monotonicity gives \(\nu_\pm(\Omega)=EX^\pm\leq E|X|<\infty\).

If \(A\in\mathcal G\) and \(P|_{\mathcal G}(A)=0\), then \(P(A)=0\), and
\(X^\pm\mathbf1_A=0\) outside \(A\), hence \(P\)-almost everywhere.
\Cref{thm:p15-nonnegative-integral-laws}, item~1, gives
\(\nu_\pm(A)=\int X^\pm\mathbf1_A\,dP=0\). So
\(\nu_\pm\ll P|_{\mathcal G}\). The measures \(\nu_\pm\) and
\(P|_{\mathcal G}\) are finite on \((\Omega,\mathcal G)\), hence
sigma-finite there. \Cref{thm:p15-radon-nikodym}, applied on the
measurable space \((\Omega,\mathcal G)\) with reference measure
\(P|_{\mathcal G}\), gives nonnegative \(\mathcal G\)-measurable
\(Z_\pm\) with values in \([0,\infty]\) and
\(\nu_\pm(A)=\int_AZ_\pm\,d(P|_{\mathcal G})\) for \(A\in\mathcal G\).

\paragraph{Item 2.}
Let \(Z_\pm\) be any nonnegative \(\mathcal G\)-measurable functions with
that property. For \(A\in\mathcal G\), item~1 of
\cref{thm:p15-w02c-p15-007-restricted-integral} gives
\[
 \int_AZ_\pm\,dP=\int_AZ_\pm\,d(P|_{\mathcal G})=\nu_\pm(A).
\]
Taking \(A=\Omega\) gives \(EZ_\pm=\nu_\pm(\Omega)=EX^\pm<\infty\).
This is the step at which \(P|_{\mathcal G}\)-integrals of the densities
become \(P\)-integrals.

\paragraph{Item 3.}
The sets \(\{Z_\pm=\infty\}=Z_\pm^{-1}(\{\infty\})\) are inverse images of
a Borel set of extended reals, so they belong to \(\mathcal G\), and so
does their union \(N\). Since \(EZ_\pm<\infty\),
\cref{thm:p15-nonnegative-integral-laws}, item~1, gives \(Z_\pm<\infty\)
\(P\)-almost everywhere. Each of the measurable sets
\(\{Z_\pm=\infty\}\) is contained in a \(P\)-null set and so has
probability zero by monotonicity. Hence \(P(N)=0\).

\paragraph{Item 4.}
Put \(Z_\pm'=Z_\pm\mathbf1_{N^c}\). These are real valued and nonnegative,
since \(Z_\pm\) is finite on \(N^c\). They are \(\mathcal G\)-measurable:
for a Borel set \(B\), \(\{Z_\pm'\in B\}\) is
\(\{Z_\pm\in B\}\cap N^c\), with \(N\) adjoined when \(0\in B\), and
\(N\in\mathcal G\). Since \(Z_\pm'\mathbf1_A=Z_\pm\mathbf1_A\) outside the
null set \(N\), \cref{thm:p15-nonnegative-integral-laws}, item~2, and
item~2 above give
\[
 \int_AZ_\pm'\,dP=\int_AZ_\pm\,dP=\nu_\pm(A)\qquad(A\in\mathcal G).
\]
In particular \(Z_+'\) and \(Z_-'\) are \(P\)-integrable.

Pointwise, \(Z=(Z_+-Z_-)\mathbf1_{N^c}=Z_+'-Z_-'\), a difference of two
real numbers at every outcome, so \(Z\) is real valued. It is
\(\mathcal G\)-measurable: for real \(c\),
\[
 \{Z<c\}=\bigcup_{q\in\mathbb Q}\bigl(\{Z_+'<q\}\cap\{Z_-'>q-c\}\bigr),
\]
because \(Z_+'<c+Z_-'\) holds exactly when some rational \(q\) satisfies
\(Z_+'<q<c+Z_-'\); every set on the right is in \(\mathcal G\).
Pointwise \(|Z|\leq Z_+'+Z_-'\), so
\(E|Z|\leq EZ_+'+EZ_-'=EX^++EX^-=E|X|<\infty\) by monotonicity and
additivity of the nonnegative integral, and \(Z\) is \(P\)-integrable.

For \(A\in\mathcal G\),
\(Z\mathbf1_A=Z_+'\mathbf1_A-Z_-'\mathbf1_A\) with nonnegative real
integrable terms. The difference rule in
\cref{thm:p15-signed-integral-laws} gives
\[
 \begin{aligned}
 \int_AZ\,dP&=\int_AZ_+'\,dP-\int_AZ_-'\,dP=\nu_+(A)-\nu_-(A)\\
            &=\int_AX^+\,dP-\int_AX^-\,dP=\int_AX\,dP,
 \end{aligned}
\]
the last equality being \cref{def:p15-integrable-functions} applied to
\(X\mathbf1_A\), whose positive and negative parts are
\(X^\pm\mathbf1_A\). Thus \(Z\) is \(\mathcal G\)-measurable and real,
its event integrals over \(\mathcal G\) are defined, and they satisfy the
identity of \cref{def:p15-general-conditional-expectation}. So \(Z\) is a
version of \(E(X\mid\mathcal G)\).

Every set used to modify a function lies in \(\mathcal G\); no subset of a
null set was adjoined to \(\mathcal G\), so no completion was used.
\end{proof}
\end{annexitem}

\begin{annexitem}{annex:p15-w02c-p15-007-sigma-finite}{thm:p15-w02c-p15-007-restriction-sigma-finite}{The trivial sigma-algebra}
\begin{proof}
\paragraph{Item 1.}
\(\mathbb R=\bigcup_{n\geq1}[-n,n]\) and \(\lambda([-n,n])=2n<\infty\), so
\(\lambda\) is sigma-finite. The sigma-algebra
\(\mathcal G=\{\emptyset,\mathbb R\}\) is contained in the Borel sets, and
\(\lambda|_{\mathcal G}(\mathbb R)=\lambda(\mathbb R)=\infty\), since
\(\lambda(\mathbb R)\geq\lambda([0,n])=n\) for every \(n\). A countable
family of members of \(\mathcal G\) whose union is \(\mathbb R\) must
contain \(\mathbb R\) itself, because a union of empty sets is empty.
That member has infinite \(\lambda|_{\mathcal G}\)-measure, so no such
cover by sets of finite measure exists.

\paragraph{Item 2.}
By \cref{thm:p15-nonnegative-integral-laws}, item~4, \(\nu\) is a measure on
the Borel sets; its values on the two members of \(\mathcal G\) are
\(\nu(\emptyset)=0\) and \(\nu(\mathbb R)=\lambda([0,1])=1\), the latter by
\cref{thm:p15-simple-integral-laws}. So \(\nu|_{\mathcal G}\) is a finite
measure. The only \(\lambda|_{\mathcal G}\)-null member of \(\mathcal G\)
is \(\emptyset\), and \(\nu(\emptyset)=0\); hence
\(\nu|_{\mathcal G}\ll\lambda|_{\mathcal G}\).

\paragraph{Item 3.}
Let \(g:\mathbb R\to[0,\infty]\) be \(\mathcal G\)-measurable, and let
\(c\) be one of its values. The set \(\{g=c\}=g^{-1}(\{c\})\) belongs to
\(\mathcal G\) and is nonempty, so it equals \(\mathbb R\). Thus \(g\) is
the constant \(c\). If \(c\) is finite, \(g=c\mathbf1_{\mathbb R}\) is a
simple function on \((\mathbb R,\mathcal G)\) and
\(\int g\,d(\lambda|_{\mathcal G})=c\cdot\lambda(\mathbb R)=c\cdot\infty\).
Under the convention of \cref{def:p15-simple-integral}, this is \(0\) for
\(c=0\) and \(\infty\) for \(c>0\). If \(c=\infty\), the simple function
\(\mathbf1_{\mathbb R}\) lies below \(g\), so by
\cref{def:p15-nonnegative-integral}
\(\int g\,d(\lambda|_{\mathcal G})\geq\lambda(\mathbb R)=\infty\). In every
case the integral lies in \(\{0,\infty\}\) and differs from
\(\nu(\mathbb R)=1\).

The hypotheses of \cref{thm:p15-radon-nikodym} other than
sigma-finiteness of the reference measure hold for
\(\nu|_{\mathcal G}\) and \(\lambda|_{\mathcal G}\), and its conclusion
fails. On the larger space the density exists: \(\mathbf1_{[0,1]}\) is
\(d\nu/d\lambda\) by construction.
\end{proof}
\end{annexitem}
